\chapter[Authentification et Gestion des Leads]{Authentification et Gestion des Leads}
\label{chap:release1}
\chapterplan{%
\chapterplanitem{}{Introduction}{sec:chap4_introduction}
\chapterplanitem{1}{Sprint 1 : Gestion des accès et des utilisateurs}{sec:sprint1}
\chapterplanitem{2}{Sprint 2 : Gestion des leads, rendez-vous et campagnes}{sec:sprint2}
\chapterplanitem{}{Conclusion}{sec:chap4_conclusion}
}

\section*{Introduction}
\phantomsection
\label{sec:chap4_introduction}

Ce chapitre constitue le socle initial de la plateforme CRM intelligente. Il regroupe les deux
premiers sprints du projet et se concentre sur les mécanismes indispensables à l'accès sécurisé
à la plateforme ainsi qu'à la gestion opérationnelle des leads et des rendez-vous.

Le Sprint 1 met en place l'authentification JWT avec un système de permissions granulaires,
la gestion des comptes utilisateurs et le contrôle d'accès basé sur les rôles (RBAC). Le
Sprint 2 introduit la gestion complète des leads, des rendez-vous et des campagnes
commerciales.

Chaque sprint est présenté selon une structure méthodologique homogène : backlog, spécification
et conception, présentation des technologies utilisées, réalisation technique, tests unitaires
et validation des résultats obtenus.

%% ══════════════════════════════════════════════════════════════════════════════
\section[Sprint 1 : Accès et utilisateurs]{Sprint 1 : Gestion des accès et des utilisateurs}
\label{sec:sprint1}
%% ────────────────────────────────────────────────────────────────────────────

\subsection{Objectifs et Backlog du Sprint 1}
\label{subsec:sprint1_backlog}

Le Sprint 1 a pour objectif de construire le socle de sécurité et de gestion des comptes
utilisateurs. Il couvre l'authentification JWT, la gestion des rôles et permissions, la
création des comptes et l'administration des accès.

\begingroup
\small
\setlength{\tabcolsep}{5pt}
\renewcommand{\arraystretch}{1.35}
\begin{longtable}{|>{\centering\arraybackslash}p{1.1cm}|>{\raggedright\arraybackslash}p{9.0cm}|>{\centering\arraybackslash}p{1.6cm}|>{\centering\arraybackslash}p{1.6cm}|}
\caption{Backlog du Sprint 1 -- Gestion des accès et des utilisateurs}\label{tab:sprint1_backlog}\\
\hline
\textbf{ID} & \textbf{User Story} & \textbf{Priorité} & \textbf{Statut} \\
\hline
\endfirsthead
\hline
\textbf{ID} & \textbf{User Story} & \textbf{Priorité} & \textbf{Statut} \\
\hline
\endhead
\hline
\endfoot
\endlastfoot
1.1 & En tant qu'\textbf{utilisateur}, je veux m'authentifier avec JWT afin d'accéder à mon espace de travail. & Élevée & Terminé \\
\hline
1.2 & En tant qu'\textbf{utilisateur}, je veux fermer ma session afin de sécuriser mon accès. & Moyenne & Terminé \\
\hline
1.3 & En tant que \textbf{Super Admin}, je veux gérer les utilisateurs (CRUD) afin d'administrer les comptes. & Élevée & Terminé \\
\hline
1.4 & En tant que \textbf{Super Admin}, je veux attribuer des rôles et permissions aux utilisateurs afin de contrôler leurs accès. & Élevée & Terminé \\
\hline
1.5 & En tant qu'\textbf{Admin RH}, je veux créer des comptes agents afin d'intégrer de nouveaux collaborateurs. & Élevée & Terminé \\
\hline
1.6 & En tant que \textbf{Super Admin}, je veux configurer les permissions granulaires afin de définir les droits d'accès. & Moyenne & Terminé \\
\hline
1.7 & En tant qu'\textbf{Admin RH}, je veux désactiver un compte agent afin de suspendre son accès. & Moyenne & Terminé \\
\hline
1.8 & En tant qu'\textbf{utilisateur}, je veux réinitialiser mon mot de passe afin de retrouver l'accès à mon compte. & Moyenne & Terminé \\
\hline
\end{longtable}
\endgroup

%% ────────────────────────────────────────────────────────────────────────────
\subsection{Spécification et Conception du Sprint 1}
\label{subsec:sprint1_spec}

\subsubsection{Diagramme de cas d'utilisation}

Le diagramme de cas d'utilisation du Sprint 1 modélise l'ensemble des interactions entre les
acteurs et le système d'authentification et de gestion des utilisateurs. Cinq acteurs sont
identifiés~: \textit{Visiteur} (non authentifié), \textit{Super Admin}, \textit{Admin RH},
\textit{Supervisor}, \textit{Quality Manager} et \textit{Agent}.

Le \textit{Visiteur} peut se connecter via son nom d'utilisateur et son mot de passe. Le
\textit{Super Admin} gère les utilisateurs, les rôles et les permissions. L'\textit{Admin RH}
crée et gère les comptes des agents. Le \textit{Supervisor} consulte les utilisateurs de son
périmètre. L'\textit{Agent} et le \textit{Quality Manager} accèdent à leurs espaces de
travail respectifs.

\subsubsection{Architecture de l'authentification}

L'authentification repose sur le standard JWT (JSON Web Token). Le flux d'authentification se
déroule comme suit~:

\begin{enumerate}
    \item L'utilisateur soumet ses identifiants (email + mot de passe) via le formulaire de
    connexion React.
    \item Le contrôleur \texttt{AuthController} reçoit la requête et appelle
    \texttt{AuthService}.
    \item Le service valide les identifiants via \texttt{PasswordHasher}, puis génère un
    token JWT contenant l'ID utilisateur, les rôles et les permissions.
    \item Le frontend stocke le token et l'inclut dans les en-têtes des requêtes suivantes.
    \item Le middleware ASP.NET Core valide le token et peuple le contexte utilisateur.
\end{enumerate}

\begin{figure}[H]
    \centering
    \fbox{\parbox{0.95\linewidth}{
    \centering\vspace{0.4cm}
    \textit{[Diagramme de séquence : Flux d'authentification JWT]}\\[0.3cm]
    Ce diagramme illustre le processus complet d'authentification~: soumission des identifiants
    depuis le frontend React, validation par AuthController et AuthService, génération du token
    JWT avec les claims (rôles, permissions), retour du token au frontend et utilisation dans
    les requêtes suivantes via le middleware de validation ASP.NET Core.
    \vspace{0.4cm}
    }}
    \caption{Diagramme de séquence -- Flux d'authentification JWT}
    \label{fig:sprint1_auth_flow}
\end{figure}

\subsubsection{Système de permissions}

Le système de permissions granulaires est implémenté via une classe \texttt{Permissions}
contenant des constantes pour chaque action du système~:

\begin{center}
\begin{tabular}{|l|l|}
\hline
\textbf{Permission} & \textbf{Description} \\
\hline
\texttt{leads.create} & Créer un nouveau lead \\
\texttt{leads.read} & Consulter les leads \\
\texttt{leads.update} & Modifier un lead \\
\texttt{leads.delete} & Supprimer un lead \\
\texttt{quality.evaluate} & Évaluer la qualité d'un appel \\
\texttt{quality.read} & Consulter les évaluations qualité \\
\texttt{attendance.read} & Consulter les pointages \\
\texttt{salary.manage} & Gérer les salaires \\
\texttt{users.manage} & Gérer les utilisateurs \\
\texttt{admin.access} & Accéder à l'administration \\
\hline
\end{tabular}
\end{center}

%% ────────────────────────────────────────────────────────────────────────────
\subsection{Présentation des Technologies Utilisées}
\label{subsec:sprint1_tech}

\subsubsection{JWT et ASP.NET Core Identity}

La plateforme utilise JWT pour l'authentification sans état~\cite{JWTRFC}. Le middleware
\texttt{Microsoft.AspNetCore.Authentication.JwtBearer} valide les tokens et peuple
automatiquement le contexte utilisateur. Le service \texttt{JwtHelper} centralise la
génération et la validation des tokens.

\subsubsection{Le filtre d'autorisation par permission}

Une classe \texttt{RequirePermissionAttribute} a été développée pour protéger les endpoints
avec des permissions spécifiques~:

\begin{itemize}
    \item Le décorateur \texttt{[RequirePermission(Permissions.Leads.Create)]} sur un contrôleur
    ou une action.
    \item Le \texttt{PermissionAuthorizationHandler} vérifie que l'utilisateur possède la
    permission requise.
    \item La configuration des permissions par rôle est définie dans la classe
    \texttt{Permissions} sous forme de mappings rôle-permissions.
\end{itemize}

%% ────────────────────────────────────────────────────────────────────────────
\subsection{Réalisation}
\label{subsec:sprint1_realisation}

\subsubsection{Interface de connexion}

La page de connexion est implémentée avec React Hook Form pour la gestion du formulaire.
Après authentification, l'utilisateur est redirigé vers son espace selon son rôle.

\begin{figure}[H]
    \centering
    \fbox{\parbox{0.9\linewidth}{
    \centering\vspace{0.4cm}
    \textit{[Capture d'écran : Page de connexion]}\\[0.3cm]
    L'interface de connexion présente un formulaire avec les champs email et mot de passe,
    un bouton de connexion et un lien de réinitialisation de mot de passe. La validation
    des champs est assurée par React Hook Form.
    \vspace{0.4cm}
    }}
    \caption{Interface de connexion}
    \label{fig:sprint1_login}
\end{figure}

\subsubsection{Gestion des utilisateurs}

L'interface d'administration des utilisateurs permet au Super Admin de créer, modifier et
désactiver des comptes. Chaque utilisateur se voit attribuer un rôle et des permissions
spécifiques.

\begin{figure}[H]
    \centering
    \fbox{\parbox{0.9\linewidth}{
    \centering\vspace{0.4cm}
    \textit{[Capture d'écran : Gestion des utilisateurs]}\\[0.3cm]
    Interface listant les utilisateurs avec leur nom, email, rôle, statut (actif/inactif)
    et les actions disponibles (modifier, désactiver, réinitialiser le mot de passe).
    \vspace{0.4cm}
    }}
    \caption{Gestion des utilisateurs}
    \label{fig:sprint1_users}
\end{figure}

\subsubsection{Tests unitaires et validation}

Les tests du Sprint 1 couvrent~: l'authentification (connexion, déconnexion, refresh token),
la gestion des utilisateurs (CRUD), le système de permissions (vérification des accès par
rôle), les validateurs FluentValidation et les scénarios d'erreur.

\subsubsection{Résultats obtenus}

Le Sprint 1 livre un socle fonctionnel pour l'authentification sécurisée et la gestion des
comptes utilisateurs. Le système de permissions granulaires permet un contrôle d'accès fin
et flexible, préparant l'extension aux fonctionnalités métier des sprints suivants. Au total,
17 contrôleurs API ont été implémentés, couvrant l'ensemble des fonctionnalités d'authentification
et de gestion des utilisateurs.

%% ══════════════════════════════════════════════════════════════════════════════
\section[Sprint 2 : Leads et rendez-vous]{Sprint 2 : Gestion des leads, rendez-vous et campagnes}
\label{sec:sprint2}

%% ────────────────────────────────────────────────────────────────────────────
\subsection{Objectifs et Backlog du Sprint 2}
\label{subsec:sprint2_backlog}

Le Sprint 2 a pour objectif de mettre en place la gestion opérationnelle des leads, des
rendez-vous et des campagnes commerciales. Il couvre la création, la qualification, le suivi
des leads, la planification des rendez-vous et la gestion des campagnes.

\begingroup
\small
\setlength{\tabcolsep}{5pt}
\renewcommand{\arraystretch}{1.35}
\begin{longtable}{|>{\centering\arraybackslash}p{1.1cm}|>{\raggedright\arraybackslash}p{9.0cm}|>{\centering\arraybackslash}p{1.6cm}|>{\centering\arraybackslash}p{1.6cm}|}
\caption{Backlog du Sprint 2 -- Leads, rendez-vous et campagnes}\label{tab:sprint2_backlog}\\
\hline
\textbf{ID} & \textbf{User Story} & \textbf{Priorité} & \textbf{Statut} \\
\hline
\endfirsthead
\hline
\textbf{ID} & \textbf{User Story} & \textbf{Priorité} & \textbf{Statut} \\
\hline
\endhead
\hline
\endfoot
\endlastfoot
2.1 & En tant qu'\textbf{Agent}, je veux créer un lead afin d'ajouter un nouveau prospect. & Élevée & Terminé \\
\hline
2.2 & En tant qu'\textbf{Agent}, je veux consulter la liste de mes leads afin de suivre mes prospects. & Élevée & Terminé \\
\hline
2.3 & En tant qu'\textbf{Agent}, je veux modifier un lead afin de mettre à jour ses informations. & Élevée & Terminé \\
\hline
2.4 & En tant qu'\textbf{Agent}, je veux qualifier un lead afin de suivre son état d'avancement. & Élevée & Terminé \\
\hline
2.5 & En tant qu'\textbf{Agent}, je veux organiser mes leads dans des folders afin de les classer. & Moyenne & Terminé \\
\hline
2.6 & En tant que \textbf{Supervisor}, je veux affecter un lead à un agent afin de distribuer la charge de travail. & Élevée & Terminé \\
\hline
2.7 & En tant qu'\textbf{Agent}, je veux planifier un rendez-vous pour un lead. & Élevée & Terminé \\
\hline
2.8 & En tant qu'\textbf{Agent}, je veux consulter mon agenda afin de visualiser mes rendez-vous. & Élevée & Terminé \\
\hline
2.9 & En tant que \textbf{Supervisor}, je veux créer une campagne afin d'organiser une action commerciale. & Élevée & Terminé \\
\hline
2.10 & En tant qu'\textbf{Agent}, je veux consulter l'historique des interactions d'un lead. & Moyenne & Terminé \\
\hline
\end{longtable}
\endgroup

%% ────────────────────────────────────────────────────────────────────────────
\subsection{Spécification et Conception du Sprint 2}
\label{subsec:sprint2_spec}

\subsubsection{Modèle de données des leads}

Le modèle de données du lead est central dans la plateforme. Il comprend les informations
suivantes~:

\begin{itemize}
    \item \textbf{Informations de contact} : nom, prénom, téléphone, email, adresse.
    \item \textbf{Qualification} : statut (HORS_CIBLE, RDV1, RDV2, RAPPEL, REFUS_COUPLE,
    REFUS_NON_INTERESSE, REFUS_CONSO, REFUS_PROJET, RDV_ANNULE, REPONDEUR).
    \item \textbf{Suivi} : date de création, date de dernier contact, agent assigné,
    campagne associée.
    \item \textbf{Organisation} : folder de classement, notes, follow-ups planifiés.
\end{itemize}

\subsubsection{Diagramme de classes — Module leads}

Le diagramme de classes du module leads est centré autour de l'entité \texttt{Lead}, qui
constitue l'élément central du CRM. Un lead est créé par un agent, peut être assigné à un
folder, rattaché à une campagne et suivi par des rendez-vous et des appels.

Les classes principales sont~:

\begin{itemize}
    \item \textbf{Lead} : entité principale avec les champs de contact, qualification et suivi.
    \item \textbf{LeadFolder} : permet de classer les leads par catégorie.
    \item \textbf{Followup} : action de suivi planifiée sur un lead.
    \item \textbf{Appointment} : rendez-vous lié à un lead et un agent.
    \item \textbf{Campaign} : campagne commerciale regroupant des leads.
\end{itemize}

\begin{figure}[H]
    \centering
    \fbox{\parbox{0.95\linewidth}{
    \centering\vspace{0.4cm}
    \textit{[Diagramme de classes du module de gestion des leads]}\\[0.3cm]
    Ce diagramme présente les entités Lead, LeadFolder, Followup, Appointment, Campaign
    et leurs relations. Un Lead appartient à un Folder et une Campaign, peut avoir plusieurs
    Followups et Appointments, et est assigné à un Agent.
    \vspace{0.4cm}
    }}
    \caption{Diagramme de classes -- Module leads}
    \label{fig:sprint2_class}
\end{figure}

\subsubsection{Diagramme de séquence — Création d'un lead}

Ce diagramme illustre le processus de création d'un lead par un agent. L'agent remplit le
formulaire de création dans l'interface React, les données sont envoyées au contrôleur
\texttt{LeadsController} via l'API REST, puis le service \texttt{LeadService} valide et
persiste le lead en base de données via le repository.

\begin{figure}[H]
    \centering
    \fbox{\parbox{0.95\linewidth}{
    \centering\vspace{0.4cm}
    \textit{[Diagramme de séquence : Création d'un lead]}\\[0.3cm]
    L'agent soumet le formulaire de création, LeadsController reçoit la requête POST,
    LeadService valide les données avec FluentValidation, le repository persiste le lead
    dans PostgreSQL, et la réponse retourne le lead créé avec son ID.
    \vspace{0.4cm}
    }}
    \caption{Diagramme de séquence -- Création d'un lead}
    \label{fig:sprint2_seq_create_lead}
\end{figure}

%% ────────────────────────────────────────────────────────────────────────────
\subsection{Présentation des Technologies Utilisées}
\label{subsec:sprint2_tech}

\subsubsection{Generic Repository et Unit of Work}

Le pattern \textbf{Generic Repository} (\texttt{IRepository<T>}) encapsule les opérations
de persistance pour chaque entité. Le \textbf{Unit of Work} (\texttt{IUnitOfWork})
coordonne les transactions entre plusieurs repositories.

\textbf{Avantages~:} Ce pattern garantit une séparation claire entre la logique métier et
la persistance, facilite les tests unitaires via le mock des repositories et assure
l'atomicité des transactions.

\subsubsection{Entity Framework Core avec PostgreSQL}

EF Core est configuré avec Npgsql pour la connexion à PostgreSQL. Les migrations sont
générées automatiquement à partir des entités du domaine. Les configurations d'entités
sont définies dans la classe \texttt{EntityConfigurations}.

\textbf{Avantages~:} EF Core permet un mapping objet-relationnel fluide, des migrations
versionnées et une exécution performante des requêtes LINQ.

%% ────────────────────────────────────────────────────────────────────────────
\subsection{Réalisation}
\label{subsec:sprint2_realisation}

\subsubsection{Interface de gestion des leads}

L'interface de gestion des leads offre une vue liste avec filtres (statut, campagne, agent,
date) et un formulaire de création/édition. L'agent peut qualifier un lead en un clic et
planifier des follow-ups.

\begin{figure}[H]
    \centering
    \fbox{\parbox{0.9\linewidth}{
    \centering\vspace{0.4cm}
    \textit{[Capture d'écran : Liste des leads]}\\[0.3cm]
    Cette capture présente la vue liste des leads avec les colonnes~: nom, téléphone, statut
    de qualification, agent assigné, campagne, date de création et actions. Des filtres
    permettent de rechercher par statut, campagne ou période.
    \vspace{0.4cm}
    }}
    \caption{Liste des leads}
    \label{fig:sprint2_leads_list}
\end{figure}

\begin{figure}[H]
    \centering
    \fbox{\parbox{0.9\linewidth}{
    \centering\vspace{0.4cm}
    \textit{[Capture d'écran : Formulaire de création d'un lead]}\\[0.3cm]
    Le formulaire de création comprend les champs~: prénom, nom, téléphone, email, adresse,
    campagne, notes initiales et statut de qualification. La validation est assurée par
    React Hook Form côté frontend et FluentValidation côté backend.
    \vspace{0.4cm}
    }}
    \caption{Création d'un lead}
    \label{fig:sprint2_lead_create}
\end{figure}

\subsubsection{Gestion des rendez-vous et de l'agenda}

L'agenda des rendez-vous permet à l'agent de visualiser et gérer ses rendez-vous par jour,
semaine ou mois. Chaque rendez-vous est lié à un lead et peut être modifié ou annulé.

\begin{figure}[H]
    \centering
    \fbox{\parbox{0.9\linewidth}{
    \centering\vspace{0.4cm}
    \textit{[Capture d'écran : Agenda des rendez-vous]}\\[0.3cm]
    Vue agenda présentant les rendez-vous de l'agent connecté avec les informations~: date,
    heure, client (lead associé), statut (confirmé, annulé, reporté) et actions disponibles.
    \vspace{0.4cm}
    }}
    \caption{Agenda des rendez-vous}
    \label{fig:sprint2_appointments}
\end{figure}

\subsubsection{Gestion des campagnes}

L'interface de gestion des campagnes permet au Supervisor de créer, suivre et analyser les
performances des actions commerciales.

\begin{figure}[H]
    \centering
    \fbox{\parbox{0.9\linewidth}{
    \centering\vspace{0.4cm}
    \textit{[Capture d'écran : Gestion des campagnes]}\\[0.3cm]
    Cette capture montre la liste des campagnes avec leur nom, date de début/fin, nombre de
    leads, taux de conversion et statut (active, terminée, en pause).
    \vspace{0.4cm}
    }}
    \caption{Gestion des campagnes}
    \label{fig:sprint2_campaigns}
\end{figure}

\subsubsection{Tests unitaires et validation}

Les tests du Sprint 2 couvrent~: la création, modification, suppression et qualification des
leads, la gestion des folders, la planification des rendez-vous, la gestion des campagnes,
ainsi que les scénarios d'erreur et les validations métier.

\subsubsection{Résultats obtenus}

Le Sprint 2 livre le cœur fonctionnel du CRM. Les agents peuvent désormais gérer leurs leads,
planifier des rendez-vous et suivre leurs campagnes. Le backend expose une API REST complète
via 17 contrôleurs, et le frontend React/TailwindCSS offre une expérience utilisateur moderne
et responsive.

%% ══════════════════════════════════════════════════════════════════════════════
\section*{Conclusion}
\phantomsection
\label{sec:chap4_conclusion}

Ce chapitre a permis de construire les fondations de la plateforme CRM. Le Sprint 1 a établi
un système d'authentification JWT robuste avec un contrôle d'accès basé sur les permissions
granulaires. Le Sprint 2 a complété ce socle par les fonctionnalités métier essentielles~:
gestion des leads, des rendez-vous et des campagnes.

À l'issue de cette première release, la plateforme dispose d'une base cohérente pour gérer
les identités, les accès et les opérations commerciales quotidiennes d'un centre d'appels.
Les tests unitaires réalisés sur l'ensemble des services renforcent la fiabilité des règles
métier critiques.
